% Generated draft: requires mathematical and source review.
\documentclass[11pt]{article}
\usepackage[margin=25mm]{geometry}
\usepackage{fontspec}
\setmainfont{Latin Modern Roman}
\usepackage{amsmath,amssymb,mathtools,unicode-math}
\setmathfont{Latin Modern Math}
\usepackage{enumitem,needspace,xurl,hyperref}
\setlength{\parindent}{0pt}
\setlength{\parskip}{5pt}
\setlength{\emergencystretch}{3em}
\newcommand{\N}{\mathbb{N}}
\newcommand{\Z}{\mathbb{Z}}
\newcommand{\Q}{\mathbb{Q}}
\newcommand{\R}{\mathbb{R}}
\newcommand{\C}{\mathbb{C}}
\newcommand{\F}{\mathbb{F}}
\newcommand{\A}{\mathbb{A}}
\newcommand{\PP}{\mathbb{P}}
\newcommand{\E}{\mathbb{E}}
\newcommand{\eps}{\varepsilon}
\newcommand{\dd}{\,\mathrm{d}}
\begin{document}
\section*{UnsolvedMath ID: 800350}
\textbf{Source problem code:} catalogue entry
\subsection{Definitions and mathematical statement}
Let \(D\) denote the parameter domain intended in the cited source, let \(B_{a,b}\) denote its \((a,b)\)-benzel for \((a,b)\in D\), and let \(I(a,b)\) denote its Conway--Lagarias invariant, with the sign convention used there. Let \(\mathcal T_{013}(B_{a,b})\) denote the set of tilings called type \(013\) in the source, described in the supplied statement as tilings by bones and left stones.

The intended question is whether
\[
\forall (a,b)\in D,\qquad I(a,b)\leq 0\ \Longrightarrow\ \mathcal T_{013}(B_{a,b})\neq\varnothing.
\]
The equality case \(I(a,b)=0\) is explicitly included.

The supplied statement also asserts the obstruction
\[
I(a,b)>0\ \Longrightarrow\ \mathcal T_{013}(B_{a,b})=\varnothing,
\]
and reports positive entries in an accompanying table whenever the invariant is nonpositive. Those observations are supplied motivation, not a proof of the proposed implication; the table itself is not supplied.

\textbf{Unresolved qualifications.} The supplied record does not specify \(D\), the geometric definition and boundary conventions of the benzel, the permitted tile shapes and orientations, the precise type \(013\) convention, or the definition and sign normalization of the invariant. The notation above therefore expresses the intended implication conditionally on those source definitions; it is not a self-contained geometric formulation.
\subsection{Short English statement}
Propp asks whether every admissible \((a,b)\)-benzel with nonpositive Conway--Lagarias invariant admits a type \(013\) tiling, described as a tiling by bones and left stones. The supplied registry labels the question open. Essential geometric definitions and parameter restrictions are missing from the supplied excerpt.
\subsection{Sources}
\begin{itemize}
\item Local UnsolvedMath registry, permanent ID \(800350\), source code \texttt{catalogue entry}, title ``Trimer Covers in the Triangular Grid --- Problem 4.'' The supplied record attributes the problem to Propp (2022) and labels it open; this status is supplied metadata, not an independently verified literature assessment.
\item Propp, \emph{Trimer Covers in the Triangular Grid: Twenty Mostly Open Problems} (2022), source item ``Problem 4 (\#2).'' Supplied reference: \url{https://arxiv.org/abs/2206.06472}. The provenance records access on 2026-10-16 and extraction method \texttt{arxiv\_latex}; these are supplied provenance details, not a claim of independent consultation.
\end{itemize}
\end{document}
